\section{Algorithms} \label{algorithms}
    \section{Basic Algorithms}

    \newpage

    \section{Bubble Sort}
        \subsection{Description}
            Bubble sort is an algorithm that sorts an unsorted list of data.
            It passes through the list, checking each pair of items to see if they need to be swapped to be in order.
            For example, if it were sorting a list of numbers into numerical order and it found the pair 2, 1, it would swap them.
            This is an example of a single bubble sort pass, with each comparison marked:
            \begin{verbatim}
                [3   5]  8   2   5   6   3
                3  [5   8]  2   5   6   3
                3   5  [8   2]  5   6   3
                3   5   2  [8   5]  6   3
                3   5   2   5  [8   6]  3
                3   5   2   5   6  [8   3]
                3   5   2   5   6   3   8
            \end{verbatim}
            Now, the 8 is in the right place, and the bubble sort will do another pass, ignoring the 8, since it must be in the right place after 1 pass, since it is the largest number.
            In general, after $n$ passes, the last $n$ items are sorted.
            Bubble sort gets its name from this property of the largets elements `bubbling' to the top (end) of the list.
            If the algorithm does a pass without swapping any items, it must have already finished sorting the list, so it will return early.

            Example of bubble sort: \\
            \begin{center}
            \begin{tabular}{c c c c c c c c}
                Pass 1: &3 &5 &8 &2 &5 &6 &3 \\
                Pass 2: &3 &5 &2 &5 &6 &3 &\textbf{8} \\
                Pass 3: &3 &2 &5 &5 &3 &\textbf{6} &\textbf{8} \\
                Pass 4: &2 &3 &5 &3 &\textbf{5} &\textbf{6} &\textbf{8} \\
                Pass 6: &2 &3 &3 &\textbf{5} &\textbf{5} &\textbf{6} &\textbf{8} \\
                Pass 7: &2 &3 &\textbf{3} &\textbf{5} &\textbf{5} &\textbf{6} &\textbf{8} \\
                \multicolumn{8}{l}{\textbf{No more swaps; sort complete.}}
            \end{tabular}
            \end{center}
            Notice that pass 6 is the same as pass 7, as only on pass 7 has the algorithm not performed any swaps.

        \subsection{Order}
            \begin{theorem}
                Bubble sort has an order of $\bigO(n)$
            \end{theorem}
            \begin{proof}
                Let $n$ be the number of elements in the list to be sorted.
                This means that (in the worst-case scenario) the bubble sort will have to do $i-1$ pass to guarantee that every item has been sorted (since after the $i$th pass, $i$ items have been sorted, $n-i$ items will be left to sort, the $i$th item must also be sorted).
                In the $i$th pass, $i-1$ items will have been sorted, so the total number of comparisons will be:
                $$\sum_{i=1}^n(n-i)=\frac{n(n-1)}{2}=\frac{1}{2}n^2+\frac{1}{2}n$$
                Taking the fastest-growing term of $\frac{1}{2}n^2+\frac{1}{2}n$ (the $\frac{1}{2}n^2$ term) and discarding the constant coefficient ($\frac{1}{2}$), this gives $n^2$, thus the order of the bubble sort algorithm is $\bigO(n^2)$.
            \end{proof}
        
        \subsection{Implementation}
            In pseudocode, this is an example implementation of a bubble sort algorithm:
            \begin{verbatim}
// Sorts a list using bubble sort; returns sorted list
// Example:
//     Input:  [3, 2, 7, 3, 5]
//     Output: [2, 3, 3, 5, 7]
//     Explanation: The unsorted input is sorted so the the numbers are returned in order
FUNC bubble_sort(list):            // Takes an unsorted list as input
    VAR n = list.UPPER_BOUND

    FOR (i = 0 TO n):                     // Goes through each pass
        VAR hasSwapped = FALSE
        FOR (j = 0 TO n - i):             // Goes through each value in the list
            IF (list[i] > list[j + 1]):   // Checks if items need swapping
                hasSwapped = TRUE
                VAR temp = list[j + 1]    // Swaps items
                list[j + 1] = list[i]
                list[i] = temp
            END IF
        END FOR

        IF (NOT hasSwapped):   // If it doesn't find any items to swap, the list must be sorted
            OUTPUT("No more swaps; sort complete.")
            RETURN list        // Return the sorted list early, since there are no more swaps
        END IF
    END FOR

    RETURN list                // Returns sorted list
END FUNC
            \end{verbatim}
    \newpage

    \section{Quick Sort}
        \subsection{Description}

        \subsection{Order}

        \subsection{Implementation}

    \newpage

    \section{Bin-Packing Algorithms}

    \newpage
    
    \section{Order of Algorithms}

    \newpage

\newpage
\section{Graphs} \label{graphs}
\newpage
\section{Algorithms on Graphs} \label{algorithms_on_graphs}
\newpage
\section{Route Inspection} \label{route_inspection}
\newpage
\section{The Travelling Salesman Problem} \label{the_travelling_salesman_problem}
\newpage
\section{Linear Programming} \label{linea_programming}
\newpage
\section{The Simplex Algorithm} \label{the_simplex_algorithm}
\newpage
\section{Critical Path Analysis} \label{critical_path_analysis}
\newpage